

This work provides quantum query and sample algorithms for implementing the Uhlmann transformation in the form of quantum circuits, building on the QSVT and the density matrix exponentiation.
We evaluate the query and sample complexities, the number of one- and two-qubit gates, and the number of qubits at any one time, for our algorithms in the purified query, the purified sample, and the mixed sample access models. 
In the purified query and sample access models, our algorithms demonstrate significant improvements in these computational costs over naive approach based on quantum state tomography and the one in Ref.~\cite{metger2023stateqipspace}.
We also consider a variant of the Uhlmann transformation for the canonical purification and propose a quantum sample algorithm in the mixed sample access model, which can be implemented with relatively few samples.
Moreover, we derive a lower bound on the query and sample complexities required for any quantum algorithm to perform the Uhlmann transformation, based on the mixedness testing in the purified query access model.



With our Uhlmann transformation algorithms, we investigate the computational costs for square root fidelity estimation.
In particular, we show that our algorithm achieves better upper bounds on query and sample complexities than the previous best algorithms.
Combined with the decoupling approach, we apply our Uhlmann transformation algorithms to two tasks: entanglement transmission and quantum state merging.
For each of these tasks, we provide a comprehensive analysis of the computational costs, demonstrating the broad applicability of the algorithms.
We also apply our Uhlmann transformation algorithm to an algorithmic implementation of the Petz recovery map.
While quantum algorithms for the Petz recovery map have been proposed using a Stinespring unitary and its inverse~\cite{gilyen2022petzmap, biswas2023noiseadapted}, we develop an algorithm that relies on the quantum channel defining the Petz recovery map, rather than its Stinespring unitary.






An important future direction is to tighten the difference between the upper and lower bounds on the computational cost for the Uhlmann transformation.
In the purified query access model, the upper and lower bounds on the query complexity obtained in this work scale differently with the rank $r_{\rm min}$; the upper bound scales linearly with $r_{\rm min}$, whereas the lower bound scales as $r_{\rm min}^{1/3}$. 
By further investigating and narrowing the discrepancy in scaling between them, we can obtain valuable insights, both theoretical and practical, into the limitations of implementing the Uhlmann transformation.


It is also worthwhile to explore how our results could be related to quantum complexity theory and quantum cryptography.
Recently, the Uhlmann transformation has been studied from a complexity-theoretic perspective, e.g., in Refs.~\cite{metger2023stateqipspace, bostanci2023unicompuhlmann, chia2024complexitytheoryquantumpromise}. 
Specifically, Ref.~\cite{metger2023stateqipspace} has provided valuable insights into complexity classes related to space and interactive proofs via the transformation, while complexity classes related to time are still not well understood.
Our Uhlmann transformation algorithm can be implemented with polynomial circuit complexity, provided that the purified state preparation unitaries are efficiently implementable and that either $1/s_{\rm min}$ or $r/\delta$ scales polynomially with the number of qubits.
This fact may offer a novel viewpoint on quantum complexity theory and quantum cryptography.


Another potential avenue is to investigate the gap between the purified and mixed sample access models.
The Uhlmann transformation is ill-defined in the mixed sample access model because it inherently requires purification.
To address this, we employed a quantum algorithm that realizes the canonical purification; however, this results in exponential costs. 
If a more efficient algorithm that allows deterministic purification rather than the random purification is developed, this exponential factor would be removed, making the Uhlmann transformation more efficiently implementable even in the mixed sample access model.
As a result, the exponential gap between the purified and mixed sample access models could be bridged.
On the other hand, if purification is easily implementable, it may break certain cryptographic assumptions, potentially leading to a ``no-go'' result for purification.
Although several results through specific cases~\cite{Chen2021ValiationalTraceFidelity, Ezzell2023MixedCompil, liu2024exponentialseparations, chen2024localtestUniInv, Sala2024SpontaneousPurif, weinstein2024efficientdetectionstrongtoweakspontaneous} are known, and efficient universal purification is impossible in the worst case~\cite{liu2025universalpurification}, a comprehensive understanding---especially whether it can be done at cost on the order of the rank---has not yet been achieved.



